% AUTO-GENERATED by scripts/gen_lean.py — section 06
\begin{theorem}[Isotropic improvement]
\label{thm:isotropic-improvement}
\lean{Kakeya4.Sec06.thm_isotropic_improvement_statement}
\uses{prop:tangent-palette}

[Isotropic improvement]
Consider the tangent data of \Cref{prop:tangent-palette}, with
\(0<k\leq1\) and \(0\leq d<1\), either in version~1, where \(j=1\), or
in version~2 with \(j=2\) and optimized narrowness \(\ell=0\).
In every substantial residual of the reference path one can, after
passing to a subsequence, find a packet \(Q=Q_s\), of thickness
\(a=N^{-s}\) and time length \(h\sim_{\log,N}a^k\), and a
time-improvement candidate of thickness
\[
 a_2=ah^{e_2},\qquad e_2>0.
\]
The candidate has incidence mass at least
\[
 \frac{h\nu_Q}{K}\left(\frac{a_2}{a}\right)^d,
 \qquad K=N^{o(1)}.
\]
Writing the candidate as \(P_{\rm new}\), its bounds
\(|P_{\rm new}|\leq Ka_2\) and
\(|\partial_wP_{\rm new}|\leq K\) hold on its assigned
trajectories throughout the same \(h\)-length block. The hidden
derivative is bounded below by \(K^{-1}\) on the counted incidences,
and its hidden curvature is at most \(K/a_2\). In version~1 the test
has the special form \(w-F(t,y)\), with \(F\) quadratic. In version~2
it is a quadratic in \((t,y,w)\).
\end{theorem}

\begin{proposition}[Physical velocity graph]
\label{a05:velocity-graph}
\lean{Kakeya4.Sec06.a05_velocity_graph_statement}

[Physical velocity graph]
For a sufficiently fine fixed power \(b\), there is a row
\(G=(G_0,G_y)\), constant on each \(W_b\) and bounded by \(K\), such
that on a substantial reference family
\begin{equation}
 |X'(t)-G(W_b)\cdot{\bf v}|\leq b^\kappa,
\end{equation}
where \(\kappa>0\) depends only on \(m\) and \(1-d\).
For every required dyadic \(s_0\geq b^\kappa\), the values of \(G\)
to width \(s_0\) have at most \(K\) possibilities in a fixed
coarse physical cell \(W_{s_0}\).
\end{proposition}

\begin{proposition}[Polynomial qualification]
\label{a05:qualification}
\lean{Kakeya4.Sec06.a05_qualification_statement}

[Polynomial qualification]
There are a subpower \(K_{\rm qual}=N^{o(1)}\) and sets \(U_N\) of
graph-reference incidences such that
\begin{equation}
 e_N:=\frac{\Lambda_{\rm ref}(U_N)}{L_{\rm ref}}
       \longrightarrow0.
\end{equation}
Every reference incidence outside \(U_N\) admits a row-valued
polynomial in \(t\), of degree at most \(M\) and coefficients
bounded by \(K_{\rm qual}\), which approximates \(G\) within
\(h^M\) on reference-labeled times of its indexed trajectory of
ordinary duration at least \(h/K_{\rm qual}\) in the same
\(Q_s\)-block, including that incidence.
\end{proposition}

\begin{lemma}[A common choice from local lists]
\label{a05:common-local-choice}
\lean{Kakeya4.Sec06.a05_common_local_choice_statement}

[A common choice from local lists]
Let \(\mu\) be an unnormalized finite measure on a finite collection of
worlds \(\omega\). In each world let \(Z(e)\in\mathbb R^D\) and
\(Y(e)\in\mathbb R^q\) be measurable functions of an incidence \(e\),
with \(|Y(e)|\leq B\). Fix an integer \(L\geq2\), bounds
\(K_1,\ldots,K_L\geq1\), and tested widths
\[
 1=s_1>s_2>\cdots>s_L=r_*,\qquad
 A=\max_{i<L}s_i/s_{i+1}.
\]
At width \(s_i\), suppose that in every half-open \(s_i\)-cell of
\(Z\) in each world, the occurring values of \(Y\) meet at most \(K_i\)
half-open \(s_i\)-cells. There is a restriction \(\mu'\leq\mu\) such
that, in every world,
\[
 \mu'(\omega)\geq
 \mu(\omega)\prod_{i=1}^L(4^DK_i)^{-1},
\]
and any two retained incidences in the same world satisfy
\[
 |Y(e)-Y(\widehat e)|
 \leq C_{D,q}(A+B)
       \bigl(|Z(e)-Z(\widehat e)|+r_*\bigr).
\]
If, within each world, both \(Z\) and \(Y\) are determined by a
discrete state, the restriction keeps or removes each whole state
of that world. No measure is renormalized.
\end{lemma}

\begin{proposition}[Conservativity of the interpolated field]
\label{a05:conservativity}
\lean{Kakeya4.Sec06.a05_conservativity_statement}

[Conservativity of the interpolated field]
In coefficient norm, or equivalently in sup norm on the required
enlarged box up to \(K\),
\[
 \|\operatorname{skew}DT\|\leq K h^{p_*-1+c_0}.
\]
\end{proposition}

\begin{lemma}[Seven-direction tensor tests]
\label{lem:tensor-kernels}
\lean{Kakeya4.Sec06.lem_tensor_kernels_statement}

[Seven-direction tensor tests]
Let \(p_i=(1,v_i)\in\R^3\), \(1\leq i\leq7\), with every triple
linearly independent.
\begin{enumerate}
\item The space of homogeneous quadratics vanishing at all \(p_i\)
has dimension at most one; every nonzero member is nonsingular.
\item The map on symmetric cubic tensors
\[
 T\longmapsto\bigl(T(p_i,p_i,\cdot)\bigr)_{i=1}^7
\]
is injective.
\item If a nonzero quadratic \(q\) vanishes at all seven points,
the kernels of the three-copy tests at ranks \(3,4,5\) are,
under the homogeneous-form correspondence,
\[
 q\,(\R^3)^*,\qquad \R q^2,\qquad \{0\},
\]
respectively.
\end{enumerate}
These assertions have polynomial conditioning. More precisely, if
\[
 |p_i|\leq A,\qquad
 |\det(p_i,p_j,p_k)|\geq A^{-1},\qquad A\geq2,
\]
the inverse of the injective map in part~2 costs at most \(CA^C\).
On the exact-conic family, with the norm, inverse norm and inverse
determinant of \(q\) also bounded by fixed powers of \(A\),
distance to each kernel in part~3 is at most \(CA^C\) times
the norm of the corresponding tests. The constants depend only
on these fixed finite-dimensional spaces and on the fixed
conditioning powers.
\end{lemma}

\begin{lemma}[Restoration of an approximate real conic]
\label{a05:real-conic}
\lean{Kakeya4.Sec06.a05_real_conic_statement}
\uses{lem:tensor-kernels}

[Restoration of an approximate real conic]
Suppose the points satisfy the conditioning bounds of
\Cref{lem:tensor-kernels}, and a real quadratic \(q(x)=x^{\rm T}
\mathsf Qx\) has
\(A^{-1}\leq\|q\|\leq A\) and \(|q(p_i)|\leq\epsilon\).
If \(\epsilon\) is smaller than a sufficiently large fixed
reciprocal power of \(A\), then
\[
 |\det\mathsf Q|\geq cA^{-C},
\]
and there are exact real zeros \(p'_i=(1,v'_i)\) of \(q\) with
\[
 |p'_i-p_i|\leq CA^C\epsilon.
\]
For sufficiently small \(\epsilon\), every triple determinant
of the shifted points has absolute value at least \((2A)^{-1}\).
\end{lemma}

\begin{proposition}[Degree restoration]
\label{prop:degree-restoration}
\lean{Kakeya4.Sec06.prop_degree_restoration_statement}
\uses{thm:isotropic-improvement}

[Degree restoration]
The potential fits and heavy intercept groups constructed above
yield a candidate with all the properties in
\Cref{thm:isotropic-improvement}, without reducing the
positive exponent \(e_2\) in \(a_2/a=h^{e_2}\).
\end{proposition}

\begin{corollary}[Zero horizon in the scalar model]
\label{a05:scalar-zero-horizon}
\lean{Kakeya4.Sec06.a05_scalar_zero_horizon_statement}

[Zero horizon in the scalar model]
For every fixed admissible choice of angular parameters and every
\(d<1\), a version~1 exact problem in \(C(d)\) has
\(H(E)=0\). In particular there is no bad problem in version~1.
\end{corollary}

